\documentclass[10pt]{article}
\usepackage[margin=20mm]{geometry}
\usepackage{amsmath,amssymb,mathtools,fontspec,unicode-math,xurl,hyperref}
\IfFontExistsTF{DejaVu Serif}{\setmainfont{DejaVu Serif}}{\setmainfont{Latin Modern Roman}}
\setmathfont{latinmodern-math.otf}
\setlength{\parindent}{0pt}\setlength{\parskip}{6pt}\sloppy
\title{Étale metaplectic covers: a review of comparison maps and categorical hypotheses}
\author{AI-assisted technical review}
\date{}
\begin{document}
\maketitle
Abstract: The specified arXiv version of Zhao’s paper contains a norm formula that does not define the asserted map of presheaves. A finite étale double cover gives an explicit failure of base-change naturality detected by a tame symbol. Three additional local statements require correction: a motivic-cohomology vanishing boundary, an omitted action in a Postnikov classification, and a global endomorphism assertion over a disconnected base. A later author-hosted manuscript removes or corrects the latter two passages, but retains the two comparison-section formulas examined here. These findings do not disprove the principal classification of étale metaplectic covers.\par
\section*{Source identity and scope}
The primary source is Yifei Zhao, Étale metaplectic covers of reductive group schemes, arXiv:2204.00610v2, 21 February 2023, 78 pages [main]. The journal identity is J. Eur. Math. Soc., DOI 10.4171/JEMS/1596 (https://doi.org/10.4171/JEMS/1596), online publication 21 January 2025. Publisher equivalence of the supplied author manuscript has not been verified. All unqualified page numbers below are one-based pages of that exact arXiv PDF, whose printed page numbers agree.\par
A separate 79-page PDF currently linked from the author’s website was retrieved [later]. The webpage labels the paper as updated in April 2026; the downloaded PDF has a distinct hash. This is a comparison source, not a replacement for the frozen primary input. The comparison here is limited to the four passages identified below, not a full collation or a claim about the publisher’s current text.\par
The detailed checks concern Sections 1.2 and 2.3, the self-cup-product and torus constructions in Sections 3–4, the strict quadratic-form and fiber-sequence structure in Section 5.1, Borel-choice dependence in Section 5.2, and dual-root and tensor-category constructions in Sections 6.1–6.4. Imported cohomology calculations, higher-categorical coherence and all applications are not independently re-proved.\par
\section*{Failure of the prescribed norm construction}
In Section 2.3.5, p. 20, the paper constructs \[\mathbb Z(n)[n]\longrightarrow \underline K_n \tag{2.16}\] from a purported presheaf map \(\mathbb Z_{\mathrm{tr}}(\mathbb G_m^{\times n})\to\underline K_n\). For an elementary finite correspondence \(Y_1\subset Y\times\mathbb G_m^n\) with coordinate units \(f_1,\ldots,f_n\), the rule takes the product of the separate norms \(\operatorname{Nm}(f_i)\) in \(K_n(Y)\).\par
Problem 1 (Major: separate norms do not give a natural map).  The rule in Section 2.3.5 fails naturality under a finite étale base change, already for \(n=2\) and smooth integral schemes over the perfect field \(\mathbb Q\).\par
\subsection*{A correspondence satisfying the stated hypotheses}
Set \[Y=\operatorname{Spec}\mathbb Q[s^{\pm1},t^{\pm1}],\qquad
 Z=\operatorname{Spec}\mathbb Q[u^{\pm1},t^{\pm1}],\qquad s=u^2.\] The map \(Z\to Y\) is finite étale of degree two. Inside \(Y\times\mathbb G_m^2\), with target coordinates \(x,y\), the equations \(x^2=s\), \(y=t\) define an integral closed correspondence finite and surjective over \(Y\), identified with \(Z\). Its two coordinate units are \(u,t\).\par
Use additive notation for \(K_2\), and write \(\{a,b\}\) for the product of two \(K_1\) units. The printed rule sends this correspondence to \[\{\operatorname{Nm}(u),\operatorname{Nm}(t)\}
       =\{-s,t^2\}.\] Pull back along \(Y'=Z\to Y\). Its image then becomes \[\{-u^2,t^2\}=2\{-u^2,t\}.\] However, the pulled-back correspondence splits as the sum of the two multiplicity-one graphs \[(u,t)\quad\text{and}\quad(-u,t).\] The factors \(x-u\) and \(x+u\) are comaximal because \(2u\) is invertible. Applying the additive rule directly after base change therefore gives \[\{u,t\}+\{-u,t\}=\{-u^2,t\}.\] The difference is \(\{-u^2,t\}\). It is nonzero: in \(\mathbb Q(u,t)\) its tame boundary at the valuation \(t=0\) is \(-u^2\neq1\) (or its inverse under the opposite convention). Thus the discrepancy is visible at the generic point and cannot disappear by Zariski sheafification. No singular correspondence, nonunit coordinate, inseparable extension or multiplicity ambiguity is involved.\par
\subsection*{Impact and bounded repair}
For this finite étale example the compatible expression is the transfer of the product: \[\operatorname{Tr}_{Z/Y}\{u,t\}
    =\{\operatorname{Nm}_{Z/Y}(u),t\}
    =\{-s,t\},\] by the projection formula. Its pullback agrees with the sum over the two graphs. Taking a norm of each coordinate separately introduces an extra degree factor.\par
The invalid presheaf construction is used to specify (2.16), then (2.19), and hence the particular comparison asserted in Theorem 2.3.8. Section 2.3.11, p. 22, invokes that theorem to define the comparison with Brylinski–Deligne extensions. A complete repair must supply a transfer-compatible motivic-to-\(K_2\) morphism on all relevant finite correspondences, its descent to the complexes, and the compatibilities subsequently used. The displayed transfer fixes this example; it is not a newly certified construction of that entire functor.\par
A standard comparison may exist by other methods. The counterexample does not say that all comparison theorems are false, nor that the main classification in Sections 4–5 is false. The Major classification concerns the explicit construction of an important comparison map, which is not defined as claimed. The same separate-norm formula remains on p. 19 of the later author-hosted PDF.\par
\section*{The boundary degree in motivic cohomology}
Problem 2 (Minor: equation (2.22) excludes a nonzero boundary group).  Page 21 attributes to Kahn’s Theorem 1.6 the assertion \[H^i(X_{\mathrm{Zar}},\mathbb Z(2))=0\qquad(i\geq4)
 \tag{2.22}\] for connected smooth \(X\) over a perfect field. The case \(i=4\) is false.\par
Take \(X=\mathbb P^2_{\mathbb Q}\). Weight-two motivic cohomology in degree four is \[H^4(X_{\mathrm{Zar}},\mathbb Z(2))\simeq
 H^2(X_{\mathrm{Zar}},\underline K_2)
 \simeq\operatorname{CH}^2(\mathbb P^2_{\mathbb Q})\simeq\mathbb Z.\] The intersection of two rational lines is a point of degree one; the degree is invariant under rational equivalence, so this class cannot vanish. This example satisfies smoothness, connectedness and perfection of the ground field.\par
The cited primary reference was actually consulted: Kahn [kahn], printed p. 401, Theorem 1.6, includes \[H^i_{\mathrm{Zar}}(X,\Gamma_{\mathrm{Zar}}(2))
   \simeq H^{i-2}_{\mathrm{Zar}}(X,\underline K_2),
       \qquad 2\leq i\leq4.\] It states zero outside the listed degrees, not at degree four. The same page identifies the codimension-two Chow term; the projective-space calculation on printed p. 411 also exhibits the nonzero \(\mathbb Z\) term.\par
The positive-degree vanishing threshold in the quoted formula should be \(i>4\). One must also distinguish global hypercohomology from cohomology sheaves. Kahn’s Remark 1.5 states that the latter occur only in degrees 1 and 2. The paper’s preceding (2.21), and its section convention in 2.3.6, concern global hypercohomology, so interpreting (2.22) as a sheaf assertion does not rescue the printed wording.\par
The next paragraph deduces the fiber of \(\mathbb Z(2)[2]\to\underline K_2\) using (2.21)–(2.22). A corrected proof can use the actual cohomology-sheaf calculation, after a valid comparison map has been supplied. The boundary correction is local and does not by itself disprove Theorem 2.3.8; neither does it repair Problem 1. The later author-hosted PDF retains the same \(i\geq4\) assertion on p. 20.\par
\section*{Missing monodromy in the preliminary classification}
Problem 3 (Minor: Lemma 1.2.7 has an unrestricted target).  Section 1.2.6 and Lemma 1.2.7, p. 13, identify \(\operatorname{Maps}_{E_1}(H,B^{n+1}A)\) with all grouplike \(E_1\) objects having \(\pi_0=H\), \(\pi_n=A\), no other homotopy groups, and a convergent Postnikov tower. The target omits the action of \(H\) on \(A\).\par
Work on the terminal site, with \(n=1\), \(H=C_2\) and \(A=C_3\), written additively. Define a strict monoidal groupoid with objects \(h\in C_2\), no arrows between distinct objects, and automorphism group \(C_3\) at each object. Put \[h\otimes k=h+k,\qquad
 a\otimes b=a+(-1)^h b\] for arrows \(a\) at \(h\) and \(b\) at \(k\). The two parenthesizations of a triple give \(a+(-1)^h b+(-1)^{h+k}c\), so associativity holds; the identity and inverse conditions are immediate. This is a grouplike, 1-truncated object with exactly the required homotopy groups. Its Postnikov tower converges.\par
Conjugation by the nonidentity component of \(C_2\) acts on \(C_3\) by inversion. In contrast, an \(E_1\) map \(C_2\to B^2C_3\) deloops to \[BC_2\longrightarrow B^3C_3.\] Its homotopy fiber is a Postnikov extension with the constant, hence trivial, \(C_2\) action on \(C_3\). The displayed nontrivial action is invariant under an equivalence preserving the identified homotopy groups. It therefore excludes our groupoid from the image, disproving essential surjectivity. Since \(C_2\) is abelian, the \(k=1\) case of Remark 1.2.8(2) is affected as well.\par
One must restrict the target to trivial action, with the appropriate identification data, or retain an action \(\rho:H\to\operatorname{Aut}(A)\) and use a twisted Postnikov class. This review proves the missing-action obstruction and the point-site counterexample, not all coherences of a replacement statement over an arbitrary site.\par
There is counterevidence to treating this as a failure of the main geometric classification. Section 4.4.1 uses the forward construction from a map to its extension, which is valid. For the torus objects used there, geometric connectedness and finite locally constant coefficients exclude the displayed component action. Moreover, the later author-hosted PDF ends Section 1.2 at 1.2.5, then starts Section 1.3 (pp. 12–13): the old 1.2.6–1.2.8 passage is absent. This is an observed version change, not an assertion about the author’s reason for it.\par
\section*{Global endomorphisms over a disconnected base}
Problem 4 (Minor: Remark 6.3.7 needs connectedness or a sheaf formulation).  Page 63 asserts \(E\simeq\operatorname{End}(\mathbf1)\) for the category of global sections of the twisted representation stack. Section 6.0.1, p. 59, has not required the base scheme to be connected.\par
Let \(S=\operatorname{Spec}\mathbb C\sqcup\operatorname{Spec}\mathbb C\), take \(G=\mathbb G_m\), trivial coefficient group \(A\), zero cover \(\mu\), and any permitted coefficient field \(E\). The global category is the product of its two untwisted component categories. Thus \[\operatorname{End}(\mathbf1)=E\times E,
 \qquad E\longrightarrow E\times E,\quad a\longmapsto(a,a).\] The idempotent \((1,0)\) is not in the image. All the stipulated reductivity and coefficient assumptions hold.\par
The endomorphism sheaf is the constant sheaf \(E\); globally one obtains \(H^0(S,E)\). For a nonempty connected \(S\) this equals \(E\). Consequently the argument that checks the claim étale locally establishes the sheaf statement, not the unrestricted global statement. Section 6.4 already imposes connectedness, so the ensuing field-oriented construction is not contradicted by this example. The later author-hosted PDF, p. 64, explicitly uses the endomorphism-sheaf formulation and adds connectedness for the global Deligne-category conclusion. This local issue is therefore already corrected in that comparison source.\par
\section*{Surviving computations and proof boundaries}
Several potentially delicate points were checked without yielding further established defects.\par
For a strict quadratic form, write \(m=\langle\check\alpha,\lambda\rangle\). The defining relation on p. 45 gives \[Q(\lambda-m\alpha)
 =Q(\lambda)+m^2Q(\alpha)-m\,b(\alpha,\lambda)=Q(\lambda).\] It implies Weyl invariance. The converse need not follow with 2-torsion coefficients: for the cocharacter lattice of \(\mathrm{SL}_2\times\mathbb G_m\) and \(A=\mathbb Z/2\), the form \(Q(m,n)=mn\) is Weyl invariant but violates the strict mixed relation. This confirms the need for the paper’s stronger condition rather than refuting it.\par
The fiber sequence of Theorem 5.1.13, p. 48, is formally compatible with the exact lattice sequence \(\Lambda_{\mathrm{sc}}\to\Lambda\to\pi_1G\) and derived Hom into \(A[2]\). Rank-one checks for \(\mathrm{SL}_2\) and \(\mathrm{PGL}_2\) preserve the torsion contributed by \(\pi_1G\); replacing derived Hom by ordinary Hom would lose it. The requisite reductive-group cohomology calculation in Proposition 5.1.11 remains an imported dependency, not independently established here.\par
The cyclic calculation behind Proposition 3.3.5 and the 2-torsion correction in Proposition 3.4.2, pp. 26,28–29, is consistent: the cocycle \(c(a,b)=ab\) on \(C_N\) contributes \(\binom N2\) after lifting a generator \(N\) times, and \[-2\binom N2\equiv N\pmod{N^2}.\] This checks the displayed cancellation, not all canonical identifications in the descent proof. The torus argument on pp. 35–37 and the rank-one 2-torsion distinctions on pp. 40–42 were examined with the latter limitation.\par
On pp. 51–52, dependence on a Borel-conjugating element is measured by \(b\), and disappears only after restriction to its radical \(\Lambda^\sharp\). The proof does not assert unconditional independence on the original torus. On pp. 60–62, if \(n_\alpha\) is the order of \(Q(\alpha)\), the strict relation implies \(n_\alpha\alpha\in\Lambda^\sharp\) and integrality of \(\check\alpha/n_\alpha\) there. The reflected lattice action is unchanged. On \(\Lambda^\sharp\), the form \(Q\) is additive and \(2Q=0\), and \(Q(n_\alpha\alpha)=0\), as needed for the central 2-torsion sign. In rank one the parity of \(n_\alpha\) gives the expected dual root datum. These arithmetic checks do not certify every descent or higher-coherence assertion.\par
The finite-coefficient choices in Section 6.4 and the applications later in the paper were not fully independently verified. No author code was executed. A bounded title/supplement/correction search identified no dedicated code package or separate erratum; the later author PDF is recorded explicitly. Search incompleteness and lack of code are not counted as defects.\par
\section*{Assessment and evidence}
Item | State\par
Overall assessment | Major: the specified motivic-to-\(K\) comparison construction fails naturality.\par
Other findings | Three bounded Minor statements with exact counterexamples.\par
Repair status | Local corrections established; the complete replacement comparison functor is not certified.\par
Principal classification | No counterexample established; whole-proof certification not claimed.\par
Later author manuscript | Two passages removed/corrected; two comparison formulas retained.\par
Scientific decision | Pending; human reviews empty.\par
\section*{Supporting checks}
The primary PDF has 996594 bytes and SHA-256\par
00000TOKEN\par
The later author PDF has 935888 bytes and SHA-256\par
00001TOKEN\par
Kahn’s reference PDF has 208934 bytes and SHA-256\par
00002TOKEN\par
The exact source-image ledger accompanies this report. Decisive primary evidence includes pp. 13,20–22,59,63; comparison images include later-author pp. 12–13,19–20,64 and Kahn PDF pp. 7,17 (printed pp. 401,411). The primary executor actually viewed these images and every final report page.\par
Three native helpers in the present conversation supplied bounded text-based checks of the torus/Postnikov statements, comparison/classification arguments, and dual-root/tensor-category constructions. The primary executor assessed their arguments, supplied and checked the primary-reference and version comparisons, wrote the report, and performed the visual inspections. Helper text reading is not recorded as visual inspection or independent acceptance. Input manifests, additional retrieval receipts, exact source/TeX/compiler/dependency hashes and byte counts, build and render exits, unchanged earlier failures, and per-page viewing receipts are retained. The controller’s UI observation is distinct from the executor’s runtime self-description; backend equivalence has not been independently established.\par
9 Y. Zhao, Étale metaplectic covers of reductive group schemes, arXiv:2204.00610v2 (https://arxiv.org/abs/2204.00610v2), 21 February 2023. Journal identity: doi:10.4171/JEMS/1596 (https://doi.org/10.4171/JEMS/1596). Exact supplied author manuscript. Y. Zhao, same title, 79-page author-hosted PDF linked from https://yifeizhao.com/, webpage update label April 2026. https://yifeizhao.com/wp-content/uploads/2026/04/part1v2.pdf. Retrieved 26 October 2026; limited passage comparison only. B. Kahn, Applications of weight-two motivic cohomology, Documenta Math. 1 (1996), 395–416, especially Remark 1.5, Theorem 1.6 and Section 5. https://webusers.imj-prg.fr/\textasciitilde{}bruno.kahn/preprints/gamma2.pdf.\par
\end{document}
